\section{Clasificación de acciones}
\label{sec:clasificacion-acciones}

Una vez finalizada la fase de identificación de jugadores (quién), continuamos con la fase de identificación de acciones (qué). Mientras que el tracking asigna identidades, equipos y posiciones a cada jugador a lo largo del vídeo, la clasificación de acciones se encarga de transformar esas trayectorias en eventos semánticos: pases, controles, robos, tiros, saques y otras situaciones narrables del partido.
Para esta tarea se ha integrado PathCRF, un modelo que formula la detección de eventos futbolísticos como una selección secuencial de aristas sobre un grafo dinámico de jugadores. En lugar de trabajar directamente sobre los frames de vídeo, PathCRF opera sobre las trayectorias $(x,y)$ de los jugadores en el campo, lo que lo hace especialmente compatible con la salida del sistema de tracking desarrollado en las fases anteriores. Esta sección describe el proceso completo: desde la adaptación de los tracks al formato requerido por el modelo hasta la generación de eventos listos para ser narrados.
\subsection{Adaptación de tracks al formato PathCRF}
El primer paso consiste en transformar la salida del tracking, almacenada como un diccionario JSON por clase y por frame, al formato tabular que espera PathCRF como entrada. Esta transformación se realiza mediante el módulo \texttt{PathCRFTracksAdapter}.
El formato de entrada de PathCRF es un \texttt{DataFrame} con columnas fijas: una columna \texttt{frame\_id}, una columna \texttt{ball\_state} y, para cada uno de los 22 jugadores (11 por equipo), 3 árbitros y el balón, un conjunto de columnas con las coordenadas $(x,y)$ proyectadas al campo y sus derivadas cinemáticas $(v_x, v_y, \textit{speed}, \textit{accel})$. Esta estructura impone dos requisitos: que cada slot tenga una identidad fija a lo largo de todo el vídeo y que todas las posiciones estén expresadas en metros sobre el terreno de juego.
\subsubsection{Asignación de slots canónicos}
El sistema de tracking produce identificadores internos (raw track IDs) que no tienen por qué mantenerse estables durante todo el partido. Para construir los slots fijos que espera PathCRF se utiliza un sistema de IDs canónicos predefinidos:
\begin{itemize}
    \item ID canónico 1: portero del equipo local (\texttt{home\_1}).
    \item IDs canónicos 3--12: jugadores de campo del equipo local (\texttt{home\_2} a \texttt{home\_11}).
    \item ID canónico 2: portero del equipo visitante (\texttt{away\_1}).
    \item IDs canónicos 13--22: jugadores de campo del equipo visitante (\texttt{away\_2} a \texttt{away\_11}).
    \item IDs canónicos 23--25: árbitros (\texttt{referee\_1} a \texttt{referee\_3}).
\end{itemize}
La asignación de tracks a slots se resuelve mediante el algoritmo húngaro. Para los jugadores de campo, se construye una matriz de costes que combina la distancia euclídea entre la posición mediana del track y una plantilla posicional de referencia, una penalización por asignar porteros a slots de campo y, cuando se utiliza asignación por frame, un coste de continuidad que favorece mantener el mismo track en el mismo slot entre frames consecutivos. Para los porteros, la selección es determinista: se escoge el track clasificado como \texttt{goalkeeper} o con rol \texttt{POR} más cercano a la portería correspondiente.
\subsubsection{Relleno de posiciones no observadas}
En muchos frames no todos los slots tienen una detección asociada. PathCRF necesita disponer de una posición para cada slot en cada instante, por lo que los huecos deben rellenarse. La estrategia consiste en dos pasos: primero se interpola linealmente entre las posiciones conocidas de cada slot, recortando los valores extrapolados a los límites del terreno de juego (105 m $\times$ 68 m). Después, cuando un slot carece por completo de observaciones durante un tramo del vídeo, se utiliza una plantilla de equipo alineada con el centroide de los jugadores visibles en ese instante. Esta plantilla sitúa a los jugadores en una formación de referencia (aproximadamente un 4-2-3-1) y se desplaza y rota para seguir la posición media del equipo.
Una estrategia similar se aplica para el balón y los árbitros. El balón se proyecta al campo cuando está disponible y se interpola en los huecos. Los árbitros se apoyan en una plantilla de tres posiciones fijas sobre el terreno de juego cuando no hay detecciones.
\subsubsection{Suavizado y cálculo de características cinemáticas}
Las trayectorias resultantes contienen ruido e irregularidades propias del tracking desde vídeo broadcast. Para reducirlo se aplica un pipeline de suavizado en dos pasadas, cada una compuesta por:
\begin{enumerate}
    \item Detección y sustitución de \textit{outliers} de movimiento aislados: si la velocidad entre un frame y el siguiente supera un umbral máximo (14 m/s para jugadores, 12 m/s para árbitros, 35 m/s para el balón) pero el puente entre el frame anterior y el posterior es plausible, el punto intermedio se corrige.
    \item Suavizado mediante mediana móvil con ventana de 9 frames, filtro Savitzky-Golay (ventana 11, orden 2) y suavizado exponencial ($\alpha = 0.35$).
    \item Limitación de la velocidad máxima por paso a un 75\% del umbral máximo.
\end{enumerate}
A partir de las trayectorias suavizadas se calculan, para cada slot y cada frame, la velocidad $(v_x, v_y)$ mediante diferencias finitas, la rapidez escalar y la aceleración. El \texttt{DataFrame} resultante contiene del orden de 26 slots $\times$ 6 columnas cinemáticas = 156 columnas, más las columnas de índice temporal.
\subsection{Modelo PathCRF}
El núcleo de la detección de acciones es PathCRF, un modelo que combina una red neuronal de atención espacio-temporal con un campo aleatorio condicional (CRF) dinámico.
\subsubsection{Arquitectura}
PathCRF modela el partido como un grafo completamente conectado que evoluciona en el tiempo. En cada instante $t$, los nodos del grafo representan a los jugadores y a nodos externos asociados a salidas del campo (\texttt{out\_left}, \texttt{out\_right}, \texttt{out\_top}, \texttt{out\_bottom}). Las aristas dirigidas entre nodos representan posibles estados de la acción: una autoarista indica que un jugador mantiene el control del balón, y una arista dirigida de un jugador $A$ hacia un jugador $B$ representa un envío de $A$ hacia $B$.
El \textit{backbone} del modelo procesa las trayectorias de los jugadores mediante bloques de atención inspirados en Set Transformer. Esto permite que el modelo capture relaciones espaciales entre jugadores sin depender del orden en que se presentan. A partir de las representaciones de los nodos, se construyen representaciones de aristas concatenando la información del nodo origen y del nodo destino.
\subsubsection{Campo aleatorio condicional}
Sobre estas representaciones, PathCRF no clasifica cada arista de forma independiente, sino que optimiza una secuencia completa de estados mediante un CRF. La puntuación total de una secuencia de aristas $y$ dada la entrada $x$ combina:
\[
s(y,x)=\sum_t e_t(y_t,x) + \sum_t a_t(y_{t-1},y_t,x)
\]
donde $e_t$ es el término de emisión (la probabilidad local de cada arista en el instante $t$) y $a_t$ es el término de transición (la coherencia entre aristas consecutivas). Además, el modelo aplica máscaras que prohíben transiciones físicamente incoherentes, como pasar de una autoarista en un jugador a una autoarista en otro jugador distinto sin pasar por una arista dirigida entre ellos.
Durante la inferencia, la secuencia más probable se obtiene mediante decodificación de Viterbi. Los eventos se extraen cuando cambia la arista seleccionada entre dos instantes consecutivos: un cambio de autoarista a arista dirigida indica el inicio de un pase, mientras que un cambio de arista dirigida a autoarista señala que el balón ha llegado a un nuevo jugador.
\subsubsection{Carga del modelo}
El modelo se carga una única vez y se mantiene en memoria mediante una caché a nivel de módulo. La carga está parametrizada por el número de \textit{trial} (en nuestro caso, el trial 120, correspondiente al checkpoint \texttt{state\_dict\_best\_acc.pt} del repositorio externo \texttt{external/pathcrf}), el dispositivo de ejecución y el fichero de checkpoint. El dispositivo se selecciona automáticamente: si hay una GPU disponible, el modelo se carga en \texttt{cuda:0}; en caso contrario, se utiliza la CPU. Esta caché evita recargar el modelo desde disco en cada ejecución, lo que es especialmente importante en el pipeline online donde se realizan múltiples pasadas de inferencia sobre ventanas consecutivas del vídeo.
El modelo original de PathCRF fue entrenado sobre el \textit{Sportec Open DFL Dataset}, que contiene datos posicionales oficiales de siete partidos de la Bundesliga alemana. Las trayectorias originales, registradas a 25 FPS, se redujeron a 5 FPS durante el entrenamiento para mejorar la eficiencia computacional, y el entrenamiento se construyó mediante ventanas deslizantes de 10 segundos sobre segmentos de juego efectivo.
\subsection{Pipeline de inferencia offline}
El pipeline completo de inferencia offline se ejecuta en una sola llamada y recorre todas las etapas: conversión de tracks, inferencia del modelo, postprocesado de edges, enriquecimiento semántico y generación de comentarios.
\subsubsection{Inferencia del modelo}
La función principal de inferencia recibe un \texttt{DataFrame} de tracking en formato PathCRF y ejecuta el modelo sobre él. Los parámetros principales son:
\begin{itemize}
    \item \texttt{use\_crf}: activa o desactiva la capa CRF. Cuando está activo, la secuencia de aristas se obtiene mediante decodificación de Viterbi sobre el grafo completo de transiciones. Cuando está desactivado, cada frame se clasifica de forma independiente (\texttt{decode = "indep"}).
    \item \texttt{window\_seconds}: ventana temporal de análisis, típicamente 10 segundos.
    \item \texttt{fps} / \texttt{sample\_freq}: frecuencia de muestreo. El modelo trabaja a 5 FPS, por lo que las trayectorias a 25 FPS se submuestrean antes de la inferencia.
\end{itemize}
La salida del modelo incluye predicciones a tres niveles: macro-previa (relaciones entre episodios), macro-posterior y micro (frame a frame). A partir de las predicciones micro se construye la secuencia de aristas, aplicando una compactación de autoaristas consecutivas para eliminar eventos espurios de muy corta duración. El umbral utilizado es de 10 frames consecutivos con la misma autoarista como mínimo para consolidar un tramo de control.

\subsection{Pipeline online con ventana deslizante}
Para que el sistema pueda generar eventos durante la ejecución del tracking, sin esperar a procesar el vídeo completo, se ha implementado un pipeline online (\textit{rolling}) que ejecuta inferencias periódicas sobre una ventana deslizante de frames.
\subsubsection{Funcionamiento}
El pipeline online mantiene un \textit{buffer} de los últimos frames procesados por el tracking. La configuración por defecto establece:
\begin{itemize}
    \item \textbf{Ventana de snapshot}: 750 frames (30 segundos a 25 FPS). Es el tamaño máximo del buffer que se conserva para cada inferencia.
    \item \textbf{Cadencia}: 40 frames (1.6 segundos). Cada cuántos frames se lanza una nueva inferencia.
    \item \textbf{Calentamiento}: 50 frames. Frames mínimos antes de la primera inferencia.
    \item \textbf{Ventana de emisión}: 40 frames, con un retardo de 50 frames respecto al frame actual. Es decir, cuando el tracking va por el frame $N$, se emiten eventos correspondientes al intervalo $[N-89, N-50]$.
\end{itemize}
Este retardo en la emisión es necesario porque PathCRF necesita ver contexto temporal alrededor de cada evento para clasificarlo correctamente. Emitir eventos correspondientes al frame actual produciría predicciones inestables, ya que el modelo no dispondría de suficiente información sobre lo que ocurre inmediatamente después.
\subsubsection{Ejecución asíncrona}
La inferencia de PathCRF es relativamente costosa (del orden de cientos de milisegundos por checkpoint), por lo que ejecutarla de forma síncrona en el bucle principal de tracking bloquearía el procesamiento de frames. Para evitarlo, el pipeline online utiliza un \texttt{ThreadPoolExecutor} con un solo \textit{worker}. Cuando se alcanza la cadencia configurada, se toma una instantánea del buffer y se envía la inferencia a ejecución en segundo plano. Mientras tanto, el bucle principal continúa acumulando frames.
Si al llegar el siguiente checkpoint la inferencia anterior aún no ha terminado, se aplica una política de descarte configurable. La política por defecto (\texttt{latest}) encola la nueva petición para ejecutarla cuando la anterior termine, mientras que la política \texttt{skip} la descarta. En la práctica, con una GPU RTX 3080, la inferencia suele completarse antes del siguiente checkpoint, por lo que rara vez se acumulan peticiones.
\subsubsection{Procesamiento de resultados}
Cuando una inferencia asíncrona termina, el hilo principal recoge el resultado de forma segura mediante un bloqueo. La secuencia de aristas predicha se proyecta sobre los IDs canónicos del sistema de tracking y los bordes correspondientes a la ventana de emisión se extraen y se pasan al módulo de postprocesado. El resultado final es un \textit{actions\_packet} que contiene: el edge activo en el frame actual, los eventos consolidados en la ventana de emisión, estadísticas del checkpoint y métricas de latencia.
\subsection{Postprocesado de edges a eventos}
Las aristas predichas por PathCRF son una secuencia densa frame a frame, pero el sistema de comentarios necesita eventos discretos y estables. La conversión de edges a eventos se realiza mediante un módulo de postprocesado diseñado para funcionar tanto en modo offline como en modo online (por bloques de emisión).
\subsubsection{Clasificación de eventos}
Cada arista $(src, dst)$ se clasifica en una de las siguientes categorías:
\begin{itemize}
    \item \textbf{Control}: cuando origen y destino son el mismo jugador (autoarista).
    \item \textbf{Pase}: cuando origen y destino son jugadores distintos.
    \item \textbf{Fuera}: cuando el origen o el destino es un nodo externo (\texttt{out\_left}, \texttt{out\_right}, \texttt{out\_top}, \texttt{out\_bottom}).
\end{itemize}
\subsubsection{Filtrado por soporte temporal}
En el modo online, cada bloque de emisión contiene edges correspondientes a una ventana de $E$ frames (40 frames en la configuración por defecto). Para que un edge se consolide como evento, debe cumplir tres condiciones:
\begin{enumerate}
    \item \textbf{Soporte mínimo}: el edge debe aparecer en al menos 5 frames del bloque.
    \item \textbf{Consecutividad}: la racha consecutiva más larga del edge debe ser de al menos 3 frames.
    \item \textbf{Ratio de soporte}: la proporción de frames con ese edge respecto al total del bloque debe ser al menos 0.15.
\end{enumerate}
Estos umbrales filtran detecciones espurias que aparecen en uno o dos frames aislados, pero permiten capturar eventos reales aunque el modelo no sea perfectamente consistente en cada frame.
\subsubsection{Histéresis y deduplicación}
Para evitar oscilaciones entre dos edges muy parecidos en bloques consecutivos, se aplica un mecanismo de histéresis. Si en el bloque anterior se emitió un edge distinto al candidato actual y el edge antiguo todavía mantiene al menos un 90\% del soporte del nuevo, se conserva el edge antiguo. Esta regla introduce una inercia que estabiliza la salida frente a pequeñas fluctuaciones del modelo.
Además, dentro de un mismo bloque de emisión solo se permite un evento (salvo que se active la opción de múltiples eventos por bloque). Esta restricción evita que el sistema genere varios eventos simultáneos que compitan por el canal de audio.
\subsection{Enriquecimiento semántico}
Una vez obtenidos los eventos base (pase, control, fuera), se aplican dos etapas de enriquecimiento semántico que añaden información específica del dominio futbolístico.
\subsubsection{Detección de tiros}
No todos los eventos de tipo \textit{kick} (pase) representan un tiro a puerta. Para distinguir los tiros reales se aplica una heurística geométrica y cinemática que evalúa cada evento de pase según varios factores:
\begin{itemize}
    \item \textbf{Distancia a portería}: el evento debe originarse a menos de 50 metros de la portería rival.
    \item \textbf{Zona de ataque}: el jugador debe encontrarse en el tercio ofensivo del campo ($x > 70$ m atacando hacia la derecha, o $x < 35$ m atacando hacia la izquierda).
    \item \textbf{Velocidad del balón}: debe superar un umbral que crece con la distancia. Para distancias cortas ($< 18$ m), el umbral es de 12 m/s; para distancias mayores, se añaden 0.25 m/s por cada metro adicional.
    \item \textbf{Dirección}: se calcula la similitud coseno entre el vector del envío y el vector hacia el centro de la portería rival. Los envíos laterales (cruces) reciben una penalización.
    \item \textbf{Bonificaciones contextuales}: se añade puntuación cuando el balón se mueve hacia la portería rival (+25), cuando el siguiente evento es una intervención del portero rival cerca de su portería (+30), o cuando aparece un evento de \textit{out} cerca de la línea de gol poco después (+25).
\end{itemize}
La puntuación total se compara con un umbral de 145 puntos. Si se supera, el evento se reclasifica como \texttt{shot} (tiro).
\subsubsection{Clasificación de saques}
El primer evento de cada episodio de juego se analiza para determinar si corresponde a un saque reglamentario:
\begin{itemize}
    \item \textbf{Córner}: la posición de inicio está a menos de 2 metros de una de las banderinas de córner.
    \item \textbf{Saque de banda}: la posición de inicio está a menos de 1 metro de una de las líneas laterales.
    \item \textbf{Saque de puerta}: la posición de inicio está dentro de un margen de 3 metros alrededor del área pequeña, sobre la línea de fondo y centrada respecto a la portería.
\end{itemize}
Los eventos clasificados como saques se etiquetan como \texttt{corner}, \texttt{throw\_in} o \texttt{goalkick} respectivamente.
\subsection{Generación de eventos narrables}
La última etapa del pipeline de acciones transforma los eventos semánticos en un formato listo para el módulo de generación de comentarios. Este proceso resuelve la identidad de los jugadores implicados, determina la posesión del balón y detecta robos.
\subsubsection{Resolución de identidad}
Cada slot canónico de PathCRF se traduce al track ID real del sistema de tracking, y a partir de él se recuperan el nombre del jugador, su posición futbolística y su equipo. Esta información procede de la fase de inferencia posicional y de la alineación introducida en la interfaz.
\subsubsection{Máquina de estados de posesión}
El sistema mantiene un estado de posesión que evoluciona con cada evento. Inicialmente la posesión está vacía. Cuando se detecta un control, la posesión se asigna al jugador que lo realiza. Cuando se detecta un pase, la posesión se transfiere al receptor. Esta máquina de estados permite filtrar eventos incoherentes, como un pase realizado por un jugador que no tiene la posesión según el histórico reciente.
\subsubsection{Detección de robos}
Un caso particularmente importante para la narración es la detección de robos de balón. Cuando el modelo predice un pase cuyo origen y destino pertenecen a equipos distintos, el evento se reclasifica como \texttt{robo}. En este caso, el jugador que recibe el balón se convierte en el actor principal de la acción y el equipo que pierde la posesión queda registrado como contexto adicional.
\subsubsection{Formato de salida}
Cada evento narrable se serializa como un objeto JSON con los siguientes campos:
\begin{itemize}
    \item \texttt{action}: tipo de acción (\texttt{pase}, \texttt{control}, \texttt{robo}, \texttt{tiro}, \texttt{corner}, \texttt{throw\_in}, \texttt{goalkick}).
    \item \texttt{event\_time\_s}: instante del evento en segundos.
    \item \texttt{player\_name} y \texttt{player\_position}: identidad del jugador que realiza la acción.
    \item \texttt{team\_name}: equipo del jugador.
    \item \texttt{action\_target}: receptor del pase, cuando existe.
    \item \texttt{field\_zone}: zona del campo (\texttt{zona de iniciación}, \texttt{zona de creación} o \texttt{zona de finalización}), calculada dividiendo el campo en tercios según la dirección de ataque del equipo.
    \item \texttt{priority}: prioridad del comentario (\texttt{high} para tiros y acciones en tercio ofensivo, \texttt{normal} para acciones en zona media, \texttt{low} para acciones en campo propio).
\end{itemize}
El resultado se almacena en un fichero \texttt{\_commentary\_events.json} que contiene dos listas paralelas: \texttt{events}, con todos los eventos detectados (incluyendo los filtrados), y \texttt{commentary\_events}, con únicamente los eventos listos para ser narrados.
\subsubsection{Integración con el pipeline principal}
En el modo online, los eventos narrables se generan bajo demanda a partir de los eventos postprocesados en cada bloque de emisión. El \textit{actions\_packet} producido por la fase de acciones incluye el último evento confirmado, que el módulo de comentarios puede consumir inmediatamente para generar y reproducir la locución correspondiente.
En resumen, la fase de clasificación de acciones completa el puente entre el tracking geométrico y la narración: convierte trayectorias de jugadores en una secuencia temporal de eventos con significado futbolístico, resuelve las identidades implicadas y produce una representación estructurada que el módulo de lenguaje puede transformar directamente en comentarios naturales.

